% ============================================================
% CONCLUSIONES
% ============================================================

\section{Conclusiones}
\label{sec:conclusiones}

Los diez ejercicios desarrollados cubren, en conjunto, el flujo completo
de procesamiento digital de señales trabajado en este curso: generación y
verificación de señales sintéticas (\texttt{WORK\_ONE}, \texttt{WORK\_TWO}),
caracterización espectral de audio real (\texttt{WORK\_THREE}),
convolución discreta (\texttt{WORK\_FOUR}), modulación de amplitud y su
comportamiento bajo distintas frecuencias de muestreo
(\texttt{WORK\_FIVE}, \texttt{WORK\_NINE}), análisis tiempo-frecuencia
mediante wavelets (\texttt{WORK\_SIX}), filtrado digital de promedio
móvil (\texttt{WORK\_SEVEN}), diseño manual de filtros IIR validado
numéricamente contra \texttt{scipy.signal} (\texttt{WORK\_EIGHT}), y
filtrado FIR/IIR combinado con la transformada de Hilbert incluyendo un
barrido de orden de filtro (\texttt{WORK\_TEN}).

Tres observaciones transversales se repiten en varias de las carpetas.
Primero, la elección de la ventana de análisis o la duración de la señal
(por ejemplo, la ventana fija de 3~s en \texttt{WORK\_THREE}, o la
extensión a 1~s del tono de 60~Hz en \texttt{WORK\_TEN} frente a los
0.05~s originales de \texttt{WORK\_FIVE}) afecta directamente la
resolución espectral obtenible por FFT, y debe documentarse junto con
cualquier resultado de frecuencia dominante. Segundo, cuando una tarea
depende de datos no disponibles en el repositorio (una tercera voz para
\texttt{WORK\_THREE} y \texttt{WORK\_SIX}, o un ambiente de Matlab para
la comparación de \texttt{WORK\_EIGHT}), el enfoque adoptado en todo el
proyecto fue documentar la ausencia explícitamente en el
\texttt{readme.md} correspondiente en lugar de fabricar el dato faltante;
este informe conserva esa misma práctica en cada subsección. Tercero, en
los ejercicios de diseño de filtros (\texttt{WORK\_SEVEN},
\texttt{WORK\_EIGHT}, \texttt{WORK\_TEN}) la validación numérica contra
una implementación de referencia (\texttt{scipy.signal}) resultó
indispensable para confirmar que una implementación manual —ya sea un
filtro de promedio móvil, un filtro IIR derivado a mano mediante la
transformada bilineal, o un filtro FIR/IIR aplicado mediante ciclo
\texttt{for}— es correcta antes de interpretar sus resultados sobre una
señal real.
